\documentclass{subfiles}
\begin{document}
    Formally, RB trees are binary trees in which the nodes store an additional bit of information: the color. 
        Hence, the informations  strored in each node are \lstinline{left, right, parent, key, color}.
        As obvious, the additional bit of information alone won't suffice our needs;
        is for this reason that Red-Black trees have to satisfy the following properties/constrains:
        \begin{enumerate}
            \item \label{RB-trees:const1} Each node is either red or black.
            \item \label{RB-trees:const2} The root of tree must be black.
            \item \label{RB-trees:const3} Each (\lstinline{NIL}) leaf is black.
            \item \label{RB-trees:const4} Each red node must have two black children.
            \item \label{RB-trees:const5} 
                For each node, any path from it to one of its descendant leaves has the same number of black nodes.
        \end{enumerate}

    For any node \(x\), property \ref{RB-trees:const5} allows us to define its black height \(\bh{x}\);
        that is, the number of black nodes in any simple path from, but not including,
        \(x\) to one of its descendant leaves.
        It follows easely that the height of the tree is the black height of the root.
        By the following result we prove that RB trees are good search trees.

    \begin{lemma}
        Let \(T\) be a red black tree with \(n\) internal nodes. 
            Then, the height of \(T\) is less or equal to \(2 \log (n + 1)\).
    \end{lemma}

    \begin{proof*}
        To prove the claim let us divide the prove in two parts: 
            we first prove that the subtree rooted at any node \(x\) has at least 
            \(2^{\bh{x}} - 1\) internal nodes; then we complete the proof by 
            showing that the height of tree is indeed at most \(2 \log(n + 1)\).

        Proceeding by induction on \(x\), let \(x\) be a node with height 0. 
            Then \(x\) is a leaf, and trivially has \(2^{\bh{x}}- 1 = 2^{0} - 1 = 0\) internal node.
            Now let \(x\) be a internal node with a positive height.
            Since \(x\) is internal node, it has two children, either or both of which are leaves.
            If the child is black, it contributes 1 to the height of \(x\) but none to its own. 
            If the child is red, it contributes neither to the parent nor itself.
            By simply noticing that the height of a child of \(x\) is less than that of \(x\),
            we can apply the inductive hypotesis and conclude that each child has at least \(2^{\bh{x} - 1} - 1\)
            internal nodes. Hence, \(x\) has at least \((2^{\bh{x} - 1} - 1) + (2^{\bh{x} - 1} - 1) + 1\) 
            internal nodes.

        To conclude the proof, note that by property \ref{RB-trees:const4},
            at least half the nodes in any simple path from, not including,
            the root to a leaf is black. Therefore, if \(h\) is the height of the tree,
            then the black height of root must be at least \(\sfrac{h}{2}\).
            By simple algebra it follows that 
            \[
                n \ge 2^{\sfrac{h}{2}}- 1 \implies h \le 2 \log (n + 1).
            \]\qed
    \end{proof*}

    As a consequence of the lemma, the operations \lstinline{INSERT, DELETE} and \lstinline{MEMBER}
    run in \(\bigO{\log n}\) time, since each can run in \(\bigO{h}\).
\end{document}
